\documentclass[a4paper,12pt,titlepage,finall]{article}

\usepackage[T1,T2A]{fontenc}
\usepackage[utf8x]{inputenc}
\usepackage[russian]{babel}
\usepackage{tikz}
\usepackage{pgfplots}
\usepackage{geometry}
\usepackage{indentfirst}
\usepackage{multirow}

\geometry{a4paper,left=30mm,top=30mm,bottom=30mm,right=30mm}
\setcounter{secnumdepth}{0}
\usepgfplotslibrary{fillbetween}

\begin{document}

\begin{titlepage}
    \begin{center}
        {\small \sc Московский государственный университет \\имени М.~В.~Ломоносова\\
        Факультет вычислительной математики и кибернетики\\}
        \vfill
        {\Large \sc Отчет по заданию №1}\\
        ~\\
        {\large \bf <<Методы сортировки>>}\\ 
        ~\\
        {\large \bf Вариант: 1/2/3/4}
    \end{center}
    \begin{flushright}
        \vfill {Выполнил:\\
        студент 104 группы\\
        Москвин~Г.~Е.\\
        ~\\
        Преподаватель:\\
        Гуляев~Д.~А.}
    \end{flushright}
    \begin{center}
        \vfill
        {\small Москва\\2026}
    \end{center}
\end{titlepage}

\tableofcontents
\newpage

\section{Постановка задачи}

В рамках данной работы решается задача экспериментального исследования эффективности двух алгоритмов сортировки массивов. Согласно варианту, исследованию подлежат метод быстрой сортировки с выбором медианы из трёх (QuickSort with median-of-three pivot) и метод сортировки вставками (Insertion Sort). Оба метода предназначены для сортировки элементов одномерного массива типа \texttt{long long} по возрастанию.

Исследование заключается в оценке вычислительной сложности алгоритмов путем подсчета двух операций:
\begin{itemize}
    \item число сравнений элементов массива;
    \item число перемещений (присваиваний/обменов) элементов в процессе сортировки.
\end{itemize}

Для объективности эксперимента тестовые массивы генерируются с разной исходной упорядоченностью. Генерация массивов выполняется отдельной функцией, которая в зависимости от параметра создаёт следующие конфигурации:
\begin{itemize}
    \item массив, упорядоченный по возрастанию (тип 1);
    \item массив, упорядоченный по убыванию (тип 2);
    \item массив со случайной расстановкой элементов (тип 3);
    \item массив с иной случайной расстановкой элементов (тип 4).
\end{itemize}

Сравнение методов сортировки производится на одних и тех же исходных массивах. Эксперименты проводятся для массивов различной длины: \\
\( n = 10, 100, 1000, 10000 \). Память под массивы выделяется динамически.

Для каждого набора параметров проводится 4 запуска, после чего результаты усредняются. Итоговые данные оформляются в виде таблиц, отражающих зависимость числа сравнений и перемещений от типа исходных данных и размера массива для каждого из исследуемых алгоритмов.

\newpage

\section{Результаты экспериментов}

В таблицах 1 и 2 приведены средние значения числа сравнений и перемещений для каждого типа начальной организации массива (типы 1–4) и для каждого размера \(n\). Значения в колонках 1–4 являются средними по 4 запускам для соответствующего типа. В последней колонке указано среднее арифметическое по четырём типам.

\begin{table}[h]
\centering
\begin{tabular}{|c|c|c|c|c|c|c|c|}
    \hline
    \multirow{2}{*}{\textbf{n}} & \multirow{2}{*}{\textbf{Параметр}} & \multicolumn{4}{|c|}{\textbf{Тип массива}} & \textbf{Среднее} \\
    \cline{3-6}
    & & \parbox{1.5cm}{\centering 1} & \parbox{1.5cm}{\centering 2} & \parbox{1.5cm}{\centering 3} & \parbox{1.5cm}{\centering 4} & \textbf{значение} \\
    \hline
    \multirow{2}{*}{10} & Сравнения & 9 & 45 & 30 & 27 & 27.75 \\
    \cline{2-7}
                        & Перемещения & 9 & 63 & 39 & 34 & 36.25 \\
    \hline
    \multirow{2}{*}{100} & Сравнения & 99 & 4950 & 2407 & 2669 & 2531.25 \\
    \cline{2-7}
                         & Перемещения & 99 & 5148 & 2508 & 2765 & 2630.00 \\
    \hline
    \multirow{2}{*}{1000} & Сравнения & 999 & 499500 & 251161 & 248806 & 250116.50 \\
    \cline{2-7}
                          & Перемещения & 999 & 501498 & 252158 & 249805 & 251115.00 \\
    \hline
    \multirow{2}{*}{10000} & Сравнения & 9999 & 49995000 & 24999865 & 25056740 & 25015401.00 \\
    \cline{2-7}
                           & Перемещения & 9999 & 50014998 & 25009868 & 25066737 & 25025400.50 \\
    \hline
\end{tabular}
\caption{Результаты работы сортировки вставками (Insertion Sort)}
\end{table}

\begin{table}[h]
\centering
\begin{tabular}{|c|c|c|c|c|c|c|c|}
    \hline
    \multirow{2}{*}{\textbf{n}} & \multirow{2}{*}{\textbf{Параметр}} & \multicolumn{4}{|c|}{\textbf{Тип массива}} & \textbf{Среднее} \\
    \cline{3-6}
    & & \parbox{1.5cm}{\centering 1} & \parbox{1.5cm}{\centering 2} & \parbox{1.5cm}{\centering 3} & \parbox{1.5cm}{\centering 4} & \textbf{значение} \\
    \hline
    \multirow{2}{*}{10} & Сравнения & 78 & 78 & 77 & 78 & 77.75 \\
    \cline{2-7}
                        & Перемещения & 45 & 78 & 67 & 70 & 65.00 \\
    \hline
    \multirow{2}{*}{100} & Сравнения & 1146 & 1376 & 1205 & 1195 & 1230.50 \\
    \cline{2-7}
                         & Перемещения & 495 & 981 & 981 & 990 & 861.75 \\
    \hline
    \multirow{2}{*}{1000} & Сравнения & 14801 & 21078 & 16353 & 16288 & 17130.00 \\
    \cline{2-7}
                          & Перемещения & 4995 & 9957 & 12230 & 12218 & 9850.00 \\
    \hline
    \multirow{2}{*}{10000} & Сравнения & 181644 & 291167 & 204205 & 202639 & 219913.75 \\
    \cline{2-7}
                           & Перемещения & 49995 & 102627 & 145406 & 145820 & 110962.00 \\
    \hline
\end{tabular}
\caption{Результаты работы быстрой сортировки (QuickSort with median-of-three)}
\end{table}

Анализ полученных результатов:

\begin{itemize}
    \item \textbf{Сортировка вставками (Insertion Sort):} 
        \begin{itemize}
            \item На уже упорядоченном массиве (тип 1) выполняется минимальное число сравнений (\(n-1\)) и перемещений (\(n-1\)), что соответствует лучшему случаю \(O(n)\).
            \item На массиве, упорядоченном по убыванию (тип 2), число сравнений максимально и равно \(n(n-1)/2\), а число перемещений также велико (\(\approx n^2/2\)), что соответствует худшему случаю \(O(n^2)\).
            \item На случайных массивах (типы 3 и 4) значения близки к среднему случаю \(O(n^2)\), но несколько меньше максимальных.
        \end{itemize}
    \item \textbf{Быстрая сортировка (QuickSort):}
        \begin{itemize}
            \item Число сравнений и перемещений на всех типах данных растёт как \(O(n \log n)\), что подтверждается экспериментальными данными (например, при \(n=10000\) значения порядка \(2\cdot10^5\)).
            \item Благодаря выбору медианы из трёх, даже на упорядоченных массивах (типы 1 и 2) алгоритм работает эффективно, избегая деградации до \(O(n^2)\).
            \item На случайных данных (типы 3 и 4) показатели близки к теоретическим оценкам.
        \end{itemize}
    \item \textbf{Сравнение при \(n=10000\):}
        \begin{itemize}
            \item Быстрая сортировка выполняет примерно в 114 раз меньше сравнений, чем сортировка вставками на обратном порядке (\(2.9\cdot10^5\) против \(5\cdot10^7\)).
            \item По перемещениям выигрыш ещё значительнее: QuickSort требует около \(1.1\cdot10^5\) перемещений, тогда как Insertion Sort — \(2.5\cdot10^7\).
            \item Таким образом, QuickSort существенно превосходит Insertion Sort при больших \(n\), что соответствует теоретическим ожиданиям.
        \end{itemize}
\end{itemize}

\newpage

\section{Структура программы и спецификация функций}

Программа реализована на языке C и состоит из трёх файлов: \texttt{main.c}, \texttt{sorts.c}, \texttt{sorts.h}. В файле \texttt{sorts.h} объявлены типы данных и прототипы функций.

\subsection{Типы данных}

\begin{verbatim}
typedef struct {
    long long comparisons;  // число сравнений
    long long moves;        // число перемещений
} SortStats;
\end{verbatim}
Структура для хранения статистики работы алгоритма сортировки.

\subsection{Функции генерации и работы с массивами}

\begin{itemize}
    \item \texttt{void generate\_array(long long *a, int n, int type)} \\
        Заполняет массив \texttt{a} длины \texttt{n} в зависимости от \texttt{type}: 
        \texttt{type=1} — упорядоченный по возрастанию, \texttt{type=2} — по убыванию, 
        \texttt{type=3} и \texttt{type=4} — случайные числа.
    \item \texttt{void copy\_array(long long *dest, long long *src, int n)} \\
        Копирует содержимое массива \texttt{src} в массив \texttt{dest}.
    \item \texttt{int is\_sorted(long long *a, int n)} \\
        Проверяет, отсортирован ли массив по возрастанию (возвращает 1, если да, иначе 0).
    \item \texttt{void print\_array(long long *a, int n, const char* message)} \\
        Выводит массив на экран (используется для отладки).
\end{itemize}

\subsection{Функции сортировки}

\begin{itemize}
    \item \texttt{void insertion\_sort(long long *a, int n, SortStats *stats)} \\
        Реализует сортировку вставками. Для каждого элемента находит его место в уже отсортированной части, сдвигая большие элементы вправо. В процессе инкрементирует поля структуры \texttt{stats}: \texttt{comparisons} при каждом сравнении, \texttt{moves} при каждом присваивании элемента.
    \item \texttt{void quick\_sort(long long *a, int n, SortStats *stats)} \\
        Обёртка для быстрой сортировки. Вызывает рекурсивную функцию \texttt{qs}.
    \item \texttt{void qs(long long *a, int left, int right, SortStats *stats)} \\
        Рекурсивная реализация быстрой сортировки. Использует функцию \texttt{partition} для разделения массива.
    \item \texttt{int partition(long long *a, int left, int right, SortStats *stats)} \\
        Выполняет разделение массива относительно опорного элемента. Опорный элемент выбирается как медиана из трёх: первого, среднего и последнего элементов. В процессе подсчитывает сравнения и обмены.
    \item \texttt{void swap(long long *a, long long *b, SortStats *stats)} \\
        Меняет местами два элемента и увеличивает счётчик перемещений на 3 (так как одно присваивание во временную переменную и два обратно).
\end{itemize}

\subsection{Основная программа (main)}

В функции \texttt{main}:
\begin{itemize}
    \item Инициализируется генератор случайных чисел.
    \item Задаются размеры массивов: 10, 100, 1000, 10000.
    \item Для каждого размера и каждого типа генерации (1..4) создаются два одинаковых массива.
    \item На первом массиве вызывается \texttt{insertion\_sort}, на втором — \texttt{quick\_sort}.
    \item Статистика каждого запуска накапливается в массивах \texttt{ins\_total} и \texttt{qs\_total}.
    \item После 4 запусков вычисляются средние значения и выводятся на экран в удобном табличном виде.
\end{itemize}

\section{Отладка программы, тестирование функций}

Отладка алгоритмов производилась на малых размерах массивов (\(n \leq 10\)) с выводом исходного и отсортированного массивов на экран. Визуально проверялось, что массив упорядочивается по возрастанию. Для автоматизации проверки использовалась функция \texttt{is\_sorted}.

Были проведены следующие тесты:
\begin{enumerate}
    \item Массив размера 5, упорядоченный по возрастанию. Обе сортировки не должны изменять порядок элементов; проверялось, что число перемещений минимально.
    \item Массив размера 5, упорядоченный по убыванию. Проверялось, что после сортировки элементы идут по возрастанию.
    \item Массив размера 10 со случайными элементами. После сортировки массив проверялся функцией \texttt{is\_sorted}.
\end{enumerate}
Все тесты прошли успешно, что подтверждает корректность реализованных алгоритмов.

\section{Анализ допущенных ошибок}

Я столкнулся с некоторыми трудностями в написании кода, а именно в написании 'sorts.h', так как раньше не сталкивался с написанием кода такого типа. Основные проблемы возникли в работе с GitHub, потому что никогда не использовал на постоянной основе.

\begin{raggedright}
\addcontentsline{toc}{section}{Список цитируемой литературы}
\begin{thebibliography}{99}
\bibitem{cs} Кормен Т., Лейзерсон Ч., Ривест Р., Штайн К. Алгоритмы: построение и анализ. Второе издание.~--- М.:<<Вильямс>>, 2005.
\end{thebibliography}
\end{raggedright}

\end{document}